% !TEX root = b-rg.tex
% Packages and macros for the reference grammar; \input by b-rg.tex.

\usepackage{graphicx} % Required for inserting images

% --- Font & Language ---
\usepackage{fontspec}
\setmainfont{Charis SIL} % Or Doulos SIL / Gentium Plus
\usepackage{polyglossia}
\setdefaultlanguage{english}

% --- Linguistic Examples ---
\usepackage{expex} % For numbered examples and glosses

\newcommand{\glphm}[3]{%
  \ex
  \begingl
  \gla #1 //
  \glb /#2/ //
  \glc #3 //
  \endgl
  \xe\ignorespaces
}

\newcommand{\glpht}[4]{%
  \ex
  \begingl
  \gla #1 //
  \glb /#2/ //
  \glb {[}#3{]} //
  \glc #4 //
  \endgl
  \xe\ignorespaces
}

% --- Glossing Abbreviations ---
% Leipzig Glossing Rules. Write grammatical glosses with the macros (\Pst,
% \Inf, \Def, \First\Sg ...), braced so expex keeps each gloss word whole:
% {swim-\Inf} {\Def=time}. Only abbreviations actually used appear in the
% list printed in frontmatter/conventions.tex (needs two XeLaTeX runs).
% Variant-to-standard mapping: planning/glossing-abbreviations.md.
\usepackage[glossaries,section]{leipzig}
\makenoidxglossaries
% Borlish-specific additions to the standard Leipzig set:
\newleipzig{agt}{agt}{agentive}
\newleipzig{cmpr}{cmpr}{comparative}
\newleipzig{coll}{coll}{collective}
\newleipzig{dsj}{dsj}{disjunctive}
\newleipzig{pot}{pot}{potential}

% --- Example Custom Command for IPA Transcription ---
\newfontfamily\ipa{Charis SIL}
\newcommand{\pht}[1]{[{\ipa #1}]}
\newcommand{\phm}[1]{/{\ipa #1}/}

\usepackage{caption}

\newcommand{\obs}{\textsuperscript{\textdagger}}
\newcommand{\bl}[1]{\emph{#1}}
\newcommand{\orth}[1]{‹~#1~›}

% --- Table of Contents ---
\setlength{\cftpartnumwidth}{2.5em} % room for part numerals like VII

% --- Drafting ---
% Placeholder for chapters that so far exist only as an outline (in comments).
\newcommand{\stub}{\par\emph{[To be written.]}\par}
